\NativeSlot{43}{36}{%
% p082_s01

\ExamQuestion{11.}{5}%
一个数，百位和百分位上都是\(6\)，其余数位上都是\(0\)，这个数是\mbox{（\hspace{1.8em}）}。\par
}
\NativeSlot{68}{42}{%
\ExamGroup{二、选择（将正确答案的序号填在括号里）。{\normalfont\kaishu\fontsize{10.8bp}{14bp}\selectfont （每题2分，共12分）}}

% p082_s02

\ExamQuestion{1.}{5}%
用简便方法计算\(602\times32\)时，下面方法正确的是\mbox{（\hspace{1.8em}）}。\par
\ExamOptions{A. \(32\times600+32\)}{B. \(32\times600+32\times2\)}{C. \(32+600+32\times2\)}{}
}
\NativeSlot{143}{36}{%
% p082_s03

\ExamQuestion{2.}{5}%
下面与\(287\)最接近的是\mbox{（\hspace{1.8em}）}。\par
\ExamOptions{A. 286.999}{B. 287.01}{C. 287.1}{}
}
\NativeSlot{193}{41}{%
% p082_s04

\ExamQuestion{3.}{5}%
一个小数的小数点先向右移动两位，又向左移动三位后是\(327.06\)，原来的小数是\mbox{（\hspace{1.8em}）}。\par
\ExamOptions{A. 32.706}{B. 327.06}{C. 3270.6}{}
}
\NativeSlot{268}{36}{%
% p082_s05

\ExamQuestion{4.}{5}%
\(24\times[623\div(12+34)-15]\)的运算顺序是\mbox{（\hspace{1.8em}）}。\par
\ExamOptions{A. 乘→除→加→减}{B. 加→除→减→乘}{C. 除→加→减→乘}{}
}
\NativeSlot{319}{55}{%
% p082_s06

\ExamQuestion{5.}{5}%
\(24\)个\(32\)相加的和除以\(12\)的\(22\)倍，商是多少？列式为\mbox{（\hspace{1.8em}）}。\par
\ExamOptions{A. \(24\times32\div12\times22\)}{B. \(24\times32\div(12\times22)\)}{C. \((24+32)\div(12\times22)\)}{}
}
\NativeSlot{418}{41}{%
% p082_s07

\ExamQuestion{6.}{5}%
在学校运动会的跳高比赛中，王兵跳了\(1.23\)米，比张宇高出了\(4\)厘米，张宇跳了\mbox{（\hspace{1.8em}）}米。\par
\ExamOptions{A. 1.27}{B. 1.53}{C. 1.19}{}
}
\NativeSlot{494}{57}{%
\ExamGroup{三、计算。{\normalfont\kaishu\fontsize{10.8bp}{14bp}\selectfont （共26分）}}

% p082_s08

\ExamQuestion{1.}{5}%
直接写得数。（8分）\par
\noindent \NativeOralCell{37\times2\times5=}\NativeOralCell{25\times101=}\NativeOralCell{20.03+6.7=}\NativeOralCell{46.25+8.6=}\par

% p082_s09

\noindent \NativeOralCell{278-32-18=}\NativeOralCell{50.8+62.7=}\NativeOralCell{182-23.6=}\NativeOralCell{55.4+9.46=}\par
}
\NativeSlot{594}{79}{%
\noindent \textbf{2.}\quad 计算下列各题，能简算的要简算。（18分）\par
\par\noindent%
\ExamTile{.323}{%
% p082_s10

\noindent \(233.2+56.4+26.8+83.6\)\par


}\hfill%
\ExamTile{.323}{%
% p082_s11

\noindent \(200-63.46-24.54\)\par


}\hfill%
\ExamTile{.323}{%
% p082_s12

\noindent \(87\times102\)\par


}\par
}
